\section{Implicit Differentiation, Gradients, Level Curves, Related Rates}

\begin{enumerate}
  \item What does it mean geometrically for a function to satisfy $F(x,y) = c$?
  \item If you move along a curve defined by $F(x,y)=c$, what is $\frac{d}{dx}F(x,y(x))$ and why?
  \item How does the chain rule lead to the formula
    \[
      \frac{\partial F}{\partial x} + \frac{\partial F}{\partial y}\frac{dy}{dx} = 0 \, ?
    \]
  \item What does the vector $(1, \frac{dy}{dx})$ represent geometrically?
  \item What is the gradient $\nabla F$, and what information does it encode?
  \item Rewrite implicit differentiation in dot product form. What does
    \[
      \nabla F \cdot (1, \tfrac{dy}{dx}) = 0
    \]
    mean geometrically?
  \item Why must the gradient be perpendicular to level curves?
  \item For the circle $F(x,y)=x^2+y^2$, what is $\nabla F$? What direction does it point?
  \item For the same circle, compute $\frac{dy}{dx}$ using implicit differentiation.
  \item Verify explicitly that the gradient is orthogonal to the tangent direction for the circle.
  \item If $x=x(t)$ and $y=y(t)$ satisfy $F(x,y)=c$, what is $\frac{dF}{dt}$ and why?
  \item Derive the related rates equation
    \[
      \frac{\partial F}{\partial x}\frac{dx}{dt} + \frac{\partial F}{\partial y}\frac{dy}{dt} = 0.
    \]
  \item What does the vector $(\frac{dx}{dt}, \frac{dy}{dt})$ represent geometrically?
  \item Rewrite related rates in dot product form. What does
    \[
      \nabla F \cdot (x', y') = 0
    \]
    mean?
  \item What is the relationship between velocity vectors and level curves?
  \item How are implicit differentiation and related rates mathematically similar?
  \item What is a directional derivative, and how does it relate to $\nabla F \cdot \mathbf{v}$?
  \item Why does $\nabla F \cdot \mathbf{v} = 0$ imply motion along a level curve?
  \item In one sentence: how are gradients, level curves, implicit differentiation, and related rates connected?
  \item Bonus: If a particle moves along a circle, why must its velocity be perpendicular to the radius?
\end{enumerate}
